% ============================================================
%  SECCIÓN — PREDIMENSIONAMIENTO DEL MURO EN MÉNSULA
% ============================================================
\section{Predimensionamiento de la estructura}
\label{sec:muro_predim}

% -----------------------------------------------------------
\subsection{Normativa y criterios aplicables}
% -----------------------------------------------------------

El predimensionamiento se realiza conforme a:

\begin{itemize}
  \item \textbf{Código Estructural} (Real Decreto 470/2021): dimensionamiento
        de estructuras de hormigón armado y definición de clases de exposición.
  \item \textbf{CTE DB-SE-C} (Seguridad Estructural — Cimientos): condiciones
        geotécnicas de la cimentación.
  \item \textbf{Eurocódigo 7 (UNE-EN 1997-1)}: diseño geotécnico.
\end{itemize}

Se adopta un muro de contención \textbf{en ménsula con zapata corrida} —
solución justificada en el apartado~\ref{sec:muro_opciones} — con los
parámetros del terreno del Sondeo~1 y la geometría que se determina a
continuación.

% -----------------------------------------------------------
\subsection{Clase de exposición y materiales}
% -----------------------------------------------------------

La presencia de sulfatos en el suelo ($3\,500$~mg~SO$_4^{2-}$/kg) y el
carácter medioambientalmente protegido de la zona condicionan directamente
la elección del hormigón.

\begin{notabox}[Clase de exposición — sulfatos]
  Según la Tabla~6.2 del Código Estructural, una concentración de
  $3\,500$~mg~SO$_4^{2-}$/kg corresponde a la clase química
  \textbf{Qa} (media agresividad). Combinada con la exposición a
  humedad y ciclos de hielo-deshielo en elementos enterrados, la clase
  de exposición resultante es \textbf{IIb+Qa}.

  Se adoptará:
  \begin{itemize}
    \item Hormigón: \textbf{HA-30/B/20/IIb+Qa}, con cemento
          \textbf{SR} (CEM~I~42{,}5~SR o CEM~III/A~42{,}5),
          resistente a sulfatos.
    \item Recubrimiento mínimo nominal (Tabla~37.2.4 CE):
          $r_{\min} = 35$~mm para clase IIb+Qa.
    \item Tolerancia de ejecución: $\Delta r = 10$~mm.
    \item \textbf{Recubrimiento nominal adoptado:}
          $r_{\text{nom}} = \resultado{50~\text{mm}}$.
    \item Hormigón de limpieza: \textbf{HM-15}, $e = 10$~cm,
          sin función estructural.
    \item Acero pasivo: \textbf{B~500~SD} ($f_{yk} = 500$~MPa).
  \end{itemize}
\end{notabox}

\begin{table}[H]
  \centering
  \caption{Materiales estructurales adoptados}
  \label{tab:materiales_muro}
  \begin{tabular}{@{}llll@{}}
    \toprule
    \textbf{Elemento} & \textbf{Hormigón} & \textbf{Cemento}
    & \textbf{Recub. nominal} \\
    \midrule
    Fuste + zapata      & HA-30/B/20/IIb+Qa & SR (CEM I 42{,}5 SR) & 50~mm \\
    Hormigón de limpieza & HM-15             & —                     & —     \\
    Acero pasivo         & B~500~SD          & —                     & —     \\
    \bottomrule
  \end{tabular}
\end{table}

% -----------------------------------------------------------
\subsection{Geometría adoptada — resumen}
\label{subsec:geom_adoptada}
% -----------------------------------------------------------

\begin{table}[H]
  \centering
  \caption{Geometría definitiva del muro en ménsula}
  \label{tab:geom_muro}
  \begin{tabular}{@{}llcl@{}}
    \toprule
    \textbf{Parámetro} & \textbf{Símbolo} & \textbf{Valor} & \textbf{Criterio} \\
    \midrule
    Altura total retenida  & $H$               & 8{,}10~m & Enunciado \\
    Altura libre del fuste & $H_f$             & 7{,}30~m & $H - h_z$ \\
    Espesor en coronación  & $e_1$             & 0{,}60~m & $H/12$ \\
    Espesor en la base     & $e_2$             & 0{,}75~m & $H/10$ \\
    Canto de la zapata     & $h_z$             & 0{,}80~m & $H/10$ \\
    Ancho total zapata     & $B$               & 5{,}60~m & Estabilidad \\
    Puntera                & $B_{\text{punta}}$ & 0{,}80~m & Estructural \\
    Talón                  & $B_{\text{talón}}$ & 4{,}05~m & $B-e_2-B_p$ \\
    Empotramiento          & $D_f$             & 1{,}00~m & Geotecnia \\
    \bottomrule
  \end{tabular}
\end{table}

% -----------------------------------------------------------
\subsection{Plano de dimensiones acotado}
\label{subsec:plano_dim}
% -----------------------------------------------------------

La figura~\ref{fig:muro_seccion} muestra la sección transversal del muro con
todas las dimensiones acotadas, el perfil del terreno y el nivel freático.

\begin{figure}[H]
\centering
\begin{tikzpicture}[scale=0.90, >=Stealth,
    dimarrow/.style={<->, >=Stealth, thin, NavyBlue!80},
    dimlabel/.style={font=\footnotesize, fill=white, inner sep=1.2pt,
                     text=NavyBlue!80},
    soil/.style={pattern=dots, pattern color=BurntOrange!50,
                 fill=BurntOrange!15},
    soilsat/.style={pattern=dots, pattern color=BurntOrange!60,
                    fill=BurntOrange!30},
    concrete/.style={fill=gray!45},
  ]

  % Origen: base zapata en (0,0). H_total = 7.30, hz=0.80,
  % B=5.00, B_tal=3.45, e2=0.75, e1=0.60, B_p=0.80
  % NF desde rasante trasdós = -2.50 m  →  y_NF = 7.30 - 2.50 = 4.80
  % Df = 1.00  →  rasante intradós en y = 1.00

  % ============================================================
  % TERRENO RETENIDO (izquierda del fuste / sobre talón)
  % ============================================================
  % Seco  (NF .. rasante trasdós)
  \fill[soil]    (-1.8, 4.80) rectangle (3.45, 7.30);
  % Saturado (base zapata .. NF)
  \fill[soilsat] (-1.8, 0.80) rectangle (3.45, 4.80);
  % Rasante trasdós
  \draw[thick, BurntOrange!70] (-1.8, 7.30) -- (3.45, 7.30);
  \fill[pattern=north east lines, pattern color=gray!60]
      (-1.8, 7.30) rectangle (3.45, 7.50);

  % ============================================================
  % TERRENO LIBRE (derecha, sobre puntera e intradós)
  % ============================================================
  \fill[soil] (4.20, 0.80) -- (5.00, 0.80) -- (5.00, 1.00)
              -- (6.80, 1.00) -- (6.80,  0.00)
              -- (5.00, 0.00) -- (4.20, 0.00) -- cycle;
  \draw[thick, BurntOrange!70] (5.00, 1.00) -- (6.80, 1.00);
  \fill[pattern=north east lines, pattern color=gray!60]
      (5.00, 1.00) rectangle (6.80, 1.15);

  % ============================================================
  % HORMIGÓN DE LIMPIEZA + ZAPATA + FUSTE
  % ============================================================
  \fill[gray!20] (-0.05,-0.12) rectangle (5.05, 0);
  \node[font=\tiny, gray!70] at (2.50,-0.06) {HM-15 ($e=10$~cm)};

  \fill[concrete] (0, 0) rectangle (5.00, 0.80);
  \draw[thick]   (0, 0) rectangle (5.00, 0.80);

  % Fuste trapezoidal: trasdós vertical x=3.45, intradós de (4.20,0.80) a (4.05,7.30)
  \fill[concrete]
    (3.45, 0.80) -- (4.20, 0.80) -- (4.05, 7.30) -- (3.45, 7.30) -- cycle;
  \draw[thick]
    (3.45, 0.80) -- (4.20, 0.80) -- (4.05, 7.30) -- (3.45, 7.30) -- cycle;

  % ============================================================
  % NIVEL FREÁTICO (símbolo estándar triángulo + rayado)
  % ============================================================
  \draw[NavyBlue!80, thick, dashed] (-1.8, 4.80) -- (3.45, 4.80);
  \node[NavyBlue!80] at (-1.5, 4.80) {$\blacktriangledown$};
  \node[NavyBlue!80, font=\scriptsize, below] at (-1.5, 4.70) {N.F.};
  \foreach \x in {-1.1, -0.7, -0.3, 0.1, 0.5, 0.9, 1.3, 1.7, 2.1, 2.5, 2.9}
    \draw[NavyBlue!60, thin] (\x, 4.80) -- (\x+0.15, 4.65);
  \node[NavyBlue!80, font=\scriptsize, right]
       at (3.50, 4.80) {$-2{,}50$~m};

  % ============================================================
  % DIAGRAMA TRIANGULAR DE EMPUJE ACTIVO (E_T)
  % ============================================================
  \fill[NavyBlue!15, opacity=0.7]
      (3.45, 0) -- (3.45, 7.30) -- (2.05, 0) -- cycle;
  \draw[NavyBlue!80, thick]
      (3.45, 7.30) -- (2.05, 0);
  % Flechas proporcionales al diagrama triangular
  \foreach \y/\len in {0.30/1.33, 1.00/1.24, 2.00/1.09, 3.00/0.94,
                       4.00/0.79, 5.00/0.64, 6.00/0.49, 7.00/0.34}
    \draw[->, NavyBlue!80, thick] (3.45-\len, \y) -- (3.45, \y);
  \node[NavyBlue!80, font=\small\bfseries] at (1.90, 4.10) {$E_T$};
  \node[NavyBlue!80, font=\tiny]           at (1.90, 3.70) {25,78 Tn/ml};

  % ============================================================
  % COTAS VERTICALES (izquierda)
  % ============================================================
  % H (total muro, desde base zapata)
  \draw[dimarrow] (-2.50, 0) -- (-2.50, 7.30);
  \node[dimlabel, rotate=90] at (-2.70, 3.65) {$H = 7{,}30$~m};

  % Hf y hz en línea intermedia
  \draw[dimarrow] (-1.95, 0.80) -- (-1.95, 7.30);
  \node[dimlabel, rotate=90] at (-2.12, 4.05) {$H_f = 6{,}50$~m};
  \draw[dimarrow] (-1.95, 0) -- (-1.95, 0.80);
  \node[dimlabel, left] at (-1.95, 0.40) {$h_z = 0{,}80$};

  % ============================================================
  % COTAS VERTICALES (derecha)
  % ============================================================
  \draw[dimarrow] (6.20, 0) -- (6.20, 1.00);
  \node[dimlabel, right] at (6.20, 0.50) {$D_f = 1{,}00$~m};

  % ============================================================
  % COTAS HORIZONTALES (abajo): tres niveles
  % ============================================================
  % Marcas verticales de referencia
  \foreach \x in {0, 3.45, 4.20, 5.00}
    \draw[thin, NavyBlue!40] (\x, -0.15) -- (\x, -1.00);

  % Nivel 1 (más bajo): B
  \draw[dimarrow] (0, -0.95) -- (5.00, -0.95);
  \node[dimlabel] at (2.50, -0.95) {$B = 5{,}00$~m};

  % Nivel 2: B_tal | e2 | Bp
  \draw[dimarrow] (0, -0.45) -- (3.45, -0.45);
  \node[dimlabel] at (1.72, -0.45) {$B_{\text{tal}} = 3{,}45$~m};
  \draw[dimarrow] (3.45, -0.45) -- (4.20, -0.45);
  \node[dimlabel, font=\scriptsize] at (3.82, -0.45) {$e_2{=}0{,}75$};
  \draw[dimarrow] (4.20, -0.45) -- (5.00, -0.45);
  \node[dimlabel, font=\scriptsize] at (4.60, -0.45) {$B_p{=}0{,}80$};

  % ============================================================
  % COTAS FUSTE (arriba)
  % ============================================================
  \draw[thin, NavyBlue!40] (3.45, 7.30) -- (3.45, 7.80);
  \draw[thin, NavyBlue!40] (4.05, 7.30) -- (4.05, 7.80);
  \draw[dimarrow] (3.45, 7.75) -- (4.05, 7.75);
  \node[dimlabel, above] at (3.75, 7.75) {$e_1 = 0{,}60$~m};

  % ============================================================
  % ETIQUETAS ESTRUCTURA
  % ============================================================
  \node[font=\bfseries\small] at (2.50, 0.40) {ZAPATA};
  \node[font=\bfseries\scriptsize, rotate=90] at (3.75, 4.10) {FUSTE};

  % Etiquetas trasdós/intradós
  \node[font=\itshape\scriptsize, NavyBlue!80, anchor=east]
       at (3.40, 2.30) {trasdós};
  \node[font=\itshape\scriptsize, BurntOrange!80, anchor=west]
       at (4.15, 2.30) {intradós};

  % Etiquetas terreno
  \node[font=\small\bfseries, BurntOrange!80]
       at (1.00, 6.15) {Terreno retenido};
  \node[font=\scriptsize, BurntOrange!80]
       at (1.00, 5.75) {(no saturado)};
  \node[font=\scriptsize\itshape, BurntOrange!80]
       at (1.00, 3.00) {(saturado)};
  \node[font=\scriptsize\itshape, BurntOrange!80]
       at (5.90, 0.50) {terreno libre};

  % Indicador de desnivel retenido (sólo informativo)
  \draw[<->, >=Stealth, thin, gray!60]
       (-0.9, 1.00) -- (-0.9, 7.30);
  \node[font=\scriptsize, gray!70, fill=white, inner sep=1pt, rotate=90]
       at (-1.05, 4.15) {Desnivel retenido: 6,30 m};

\end{tikzpicture}
\caption{Sección transversal tipo del muro en ménsula
         (dimensiones en m; N.F.\ a $-2{,}50$~m desde rasante trasdós;
         empuje activo total $E_T = 25{,}78$~Tn/ml)}
\label{fig:muro_seccion}
\end{figure}

% -----------------------------------------------------------
\subsection{Detalles constructivos}
\label{subsec:detalles_constructivos}
% -----------------------------------------------------------

\subsubsection{Drenaje del trasdós}

La presencia del nivel freático a $-2{,}50$~m hace imprescindible un sistema
de drenaje en el trasdós del muro para evitar sobrepresiones hidrostáticas no
previstas en el diseño. Se dispondrá:

\begin{enumerate}
  \item \textbf{Tubo drenante} Ø~100~mm perforado, envuelto en geotextil no
        tejido (150~g/m²), colocado en la base del trasdós sobre la zapata.
  \item \textbf{Capa drenante} de grava limpia (árido 20/40~mm), $e \geq 0{,}30$~m,
        entre el trasdós del hormigón y el terreno retenido.
  \item \textbf{Geotextil} de separación entre la grava y el relleno, para
        evitar la colmatación.
  \item \textbf{Lámina impermeabilizante} (HDPE, $e \geq 1{,}5$~mm) sobre el
        trasdós del fuste, adherida con imprimación bituminosa.
\end{enumerate}

\subsubsection{Juntas de dilatación}

Dado que el muro presenta una longitud total de $L = 150$~ml, se proyectan 
juntas de dilatación cada $15$~m con objeto de permitir los movimientos de 
expansión y contracción del hormigón producidos por las variaciones térmicas, 
reduciendo así los esfuerzos internos y el riesgo de fisuración de la estructura.

De este modo, el muro queda dividido en \textbf{10~tramos} de $15$~m de longitud.

Las juntas se ejecutarán mediante material compresible y sellado elástico para 
garantizar la estanqueidad y permitir la deformación relativa entre tramos. Asimismo, 
en las juntas se interrumpirá la continuidad de determinadas armaduras longitudinales, 
disponiendo elementos adecuados para asegurar la transmisión de esfuerzos cortantes 
entre paños.

\subsubsection{Hormigón de limpieza}

Antes del ferrallado de la zapata se extenderá una capa de
\textbf{HM-15 de 10~cm de espesor} como superficie de trabajo limpia y
nivelada. El HM-15 no se incluye en los cálculos estructurales.

\subsubsection{Ejecución por tongadas}

El hormigonado del fuste se realizará por \textbf{tongadas de 50~cm de
altura máxima}, mediante encofrado metálico a dos caras, con vibrado
interno de cada tongada antes de verter la siguiente.


